\section{本章小结}
\label{sec:lie-summary}

本章开头问：怎样说清一个计算“不受旋转影响”还是“随旋转同步变化”，以及怎样在只有三个自由度的旋转上做计算而不破坏约束。第一个问题的回答分两层。群只规定变换怎样复合：结合律、单位元和逆元保证变换可以接连施加和撤销，复合的次序一般不能交换。子群是对运算封闭的允许变换，而不是任意子集；同态保持复合，它的核记录丢失了哪些区分，像记录实际表示了哪些元素，元素个数相同的两个群未必同构。

群作用再规定每个变换怎样改变一个具体对象。同一个群可以作用在点、点云和图像上，公式各不相同，都要核对单位元与复合两条规则。轨道收集通过变换可以互相到达的对象，稳定子收集不改变给定对象的变换，二者分别是对象集合与群的子集，由陪集一一联系。在输入与输出的作用都给定后，不变映射在每条轨道上取常值，等变映射的输出按指定规则随输入变化；重心对置换不变、对旋转等变，距离矩阵对旋转不变、对置换在共轭作用下等变。不变是输出作用平凡时的等变，不变性也不保证保留了区分轨道所需的信息。

第二个问题由矩阵Lie群回答。$\mathrm{GL}(n,\mathbb R)$是开集，$\mathrm O(n)$与$\mathrm{SO}(n)$是正则水平集，乘法与求逆都光滑，左平移把单位元处的切空间搬到每一点。单位元处的切空间即Lie代数，可以做线性组合，交换子记录变换的不交换；它的元素是速度，不是群中的变换。矩阵指数由幂级数定义，把Lie代数方向变成单参数子群，并在单位元附近给出局部坐标，但一般既不单射也不满射。它只依赖群结构，与依赖度量的黎曼指数是两回事，只在$\mathrm{SO}(n)$配Frobenius度量这类特殊情形中吻合。

三维旋转没有全局无奇异的三参数坐标，这由$\mathrm{SO}(3)$的紧性直接推出。旋转矩阵直接作用于向量；轴角由Rodrigues公式给出，在半圈旋转与恒等旋转处有冗余，旋转向量在模长小于$\pi$的开球内是局部坐标；Euler角在万向节锁处秩下降，那是坐标的奇异而非旋转的退化。单位四元数以乘法完成复合，在每一点附近都给出正常坐标，代价是$q$与$-q$表示同一旋转，这一双覆盖不是额外的物理自由度。刚体位姿把旋转与平移合成$\mathrm{SE}(3)$，复合时前一次的平移要随后一次的旋转转动，方向不受平移影响；同一个齐次矩阵既可读作主动变换，也可读作坐标系换算。

最后，多种记录代表同一对象时，用群作用把同一轨道合并。商集合与商拓扑总能定义，但光滑结构需要条件：自由且适当的光滑作用给出维数为$\dim M-\dim G$的商流形，闭子群给出齐性空间，例如$\mathrm{SO}(3)/\mathrm{SO}(2)\cong S^2$；$\mathrm{SO}(2)$作用于平面的商有端点，无理旋转的商不是Hausdorff空间，都说明这些条件不能省略。射影空间合并非零倍数，用仿射图记录点、用可逆矩阵诱导射影变换；Grassmann空间合并张成同一子空间的不同基，标架比子空间多出的自由度可以用投影矩阵或水平方向消去。由点对应实际估计刚体变换与射影变换，在第~\ref{chap:spatial-map-estimation}章研究；正交矩阵与子空间上的优化算法留给卷五。
